\section{Discussion and Threats to Validity}
\label{sec:threats}

\subsection{Engineering Takeaways for High-Velocity CI}
Our empirical results highlight three practical lessons for teams scaling merge queues in the era of AI coding agents:
\begin{enumerate}
    \item \textbf{Prevent Conflicts Before Speculating}: Speculative linear trains waste nearly one-third of CI cycles (\SpecLinInvalidRatePct{}\%) under $5\times$ arrival surges. Reordering the ready queue into graph-disjoint fast lanes eliminates intra-batch semantic conflicts proactively and delivers the single largest improvement in latency and compute.
    \item \textbf{Guard Against Blind Flaky Retries vs.\ Clean-PR False Rejections}: Unconditional retry-on-failure (\texttt{FlakeAwareBatch}) severely degrades throughput and compute (\FlakeBatchCICost{}~min/PR). Conversely, disabling retries entirely (\texttt{SSMQ-NoFlaky}) lowers survivor-only latency (\AblNoFlakyLat{}~min) at the cost of falsely rejecting additional clean PRs that hit transient flakes (e.g., on \texttt{kubernetes} seed 1001, \texttt{SSMQ} merges 139 PRs vs.\ 136 for \texttt{SSMQ-NoFlaky}). Adaptive posterior gating balances clean-PR acceptance against retry burn.
    \item \textbf{Fail-Safe Agentic Integration}: Embedding simulated causal fault attribution inside a deterministic binary-bisection fallback yields consistent fault-isolation gains (\SSMQBlameRounds{} vs.\ \StaticBlameRounds{} rounds against static FIFO bisection) without allowing misattributions to bypass CI verification.
\end{enumerate}

\subsection{Threats to Validity}
\textbf{Construct Validity.} First, while PR file sets, churn, and co-change graphs come from 600 real merged GitHub PRs (with GitHub API file lists capped at 40 files/PR and merge timestamps used as arrival proxies), individual defects and pairwise conflicts are synthetically injected. Crucially, the conflict generator uses the same warm-up impact-cone overlap predicate that SSMQ's graph-disjoint batcher avoids, so our current replay isolates the scheduling ceiling under accurate impact graphs rather than off-graph latent conflicts. Second, the simulator checks injected faults within each verified job rather than compiling a cumulative mainline tree, and causal attribution is modeled via parameterized Top-1/Top-2 accuracy rather than live LLM execution on stack traces. Third, our \texttt{UberSubmitQueue} baseline is a restricted linear-prefix risk-budgeted proxy (fixed $0.30$ risk budget and single root-bisection hedge) rather than Uber's full outcome-tree value planner~\cite{ananthanarayanan2019keeping}. Addressing all three via off-graph conflict controls and a live shadow merge queue is scheduled as our Gate 3 next-run roadmap.

\textbf{Internal & Conclusion Validity.} Random simulation variance is assessed across 20 seeds per repository ($60$ pooled records per configuration across 3 fixed repos) with an A/A seed-block check (\AANoiseThroughputPct{}\% throughput diff, \AANoiseLatencyPct{}\% P95 diff, \AANoiseComputePct{}\% compute diff). Because seeds are reused across repositories and P95 latency conditions on merged PRs, marginal Cliff's $\delta$ values are descriptive effect summaries rather than formal multi-repo hypothesis tests; paired per-seed inferential contrasts are slated for Gate 3 follow-up approval.

\textbf{External Validity.} Our traces span three large open-source ecosystems (Go, TypeScript, JavaScript/Rust monorepo). Repositories with radically different CI build-time variance, monorepo build-cache sharing, or multi-hour hardware-in-the-loop suites may exhibit different optimal batch caps $K_{\max}$ and risk thresholds $\tau_{\text{risk}}$.
